% Source: https://artofproblemsolving.com/community/c3h1749495_length_of_segment
% Stable identity: aops-topic-1749495-q01
Let $P$ be a corner of the box and $Q$ the opposite corner. Walk from $P$ along the bottom face to the corner below $Q$: this is a diagonal of an $a\times b$ rectangle, of length $e=\sqrt{a^2+b^2}$ by the Pythagorean theorem. That diagonal is perpendicular to the vertical edge of length $c$ that ends at $Q$, so by the Pythagorean theorem again
\[
PQ=\sqrt{e^2+c^2}=\boxed{\sqrt{a^2+b^2+c^2}}.
\]

\begin{asy}
usepackage("amsmath");
size(6cm);
defaultpen(fontsize(10pt));
real a=4,c=2.4,bb=2.2;pair v=bb*(0.55,0.4);
pair P0=(0,0),P1=(a,0),P2=(a,c),P3=(0,c);
pair Q0=P0+v,Q1=P1+v,Q2=P2+v,Q3=P3+v;
draw(P0--P1--P2--P3--cycle,linewidth(0.9));
draw(P1--Q1--Q2--Q3--P3,linewidth(0.9));draw(P2--Q2,linewidth(0.9));
draw(P0--Q0,dashed);draw(Q0--Q1,dashed);draw(Q0--Q3,dashed);
draw(P0--Q1,blue+dashed+linewidth(0.8));
draw(P0--Q2,red+linewidth(1.1));
dot(P0);dot(Q2);
label("$P$",P0,SW);label("$Q$",Q2,NE);
label("$a$",(P0+P1)/2,S);label("$b$",(P1+Q1)/2,SE);label("$c$",(Q1+Q2)/2,E);
label("$e$",P0+0.6*(Q1-P0),S,blue);
label("\mbox{}",(min(currentpicture,user=true)+max(currentpicture,user=true))/2,black); // keeps the text after the figure black when LaTeX typesets the labels inline
\end{asy}

No two vertices are farther apart: two vertices differ in $0$, $1$, $2$ or $3$ of the coordinates (in a coordinate system along the edges), giving distances $0$, an edge, a face diagonal or the space diagonal, and the space diagonal is the largest. For a cube of edge $a$ the space diagonal is $\sqrt{3a^2}=a\sqrt3$.
